\subsection{The remaining character and the limit of a Pell-unit argument}
\label{sec:kt-third-unit-character}

This subsection tests a specific proposed mechanism for the third unit in
KT-EXPLICIT-TWO-UNIT-FAMILY.  It identifies the character that must be
controlled, proves that units from multiquadratic fields have zero
projection to that character, and shows that requiring the entire
cyclotomic compositum to be $p$-rational is too strong.  None of these
statements asserts failure of the third unit itself.

\begin{lemma}[KT-PADIC-BRANCH-RESIDUE]
Let $p>2$ be prime and let $\psi$ be a nontrivial even Dirichlet
character whose order divides $p-1$, with values in $\mathbf Z_p$.
For all $s,t\in\mathbf Z_p$,
\[
 L_p(s,\psi)\in\mathbf Z_p,\qquad
 L_p(s,\psi)\equiv L_p(t,\psi)\pmod p.                 \tag{TU.1}
\]
There is no restriction that the conductor be prime to $p$.
\end{lemma}
\begin{proof}
Write $f$ for the conductor.  Choose an integer $c>1$, prime to $pf$,
with $\psi(c)\ne1$.  Such a residue class exists because $\psi$ is
nontrivial; imposing that it be prime to $p$ adds no obstruction.
Since the order of $\psi$ divides $p-1$, reduction is injective on its
values, so $1-\psi(c)$ is a unit.

Knospe--Washington, Corollary~2.5 on author-PDF p.~5, gives
\[
 -(1-\psi(c)\langle c\rangle^{1-s})L_p(s,\psi)
 =\lim_{n\to\infty}\sum_{\substack{1\le a\le fp^n\\p\nmid a}}
     \psi\omega^{-1}(a)\langle a\rangle^{-s}
     \epsilon_{a,c,fp^n},                                  \tag{TU.2}
\]
where
$\epsilon_{a,c,F}=(c-1)/2-\{-aF^{-1}\}_c\in\tfrac12\mathbf Z$.
Every summand is integral.  For $s,t\in\mathbf Z_p$,
$\langle a\rangle^{-s}\equiv\langle a\rangle^{-t}\equiv1\pmod p$.
Thus each finite sum is integral and its difference between $s$ and
$t$ is in $p\mathbf Z_p$.  Taking the two convergent limits preserves
these properties.  The factor on the left is a unit, congruent to
$1-\psi(c)$ independently of $s$.  Division proves (TU.1).
The cited domain $|s|<p^{1-1/(p-1)}$ contains $\mathbf Z_p$.
\end{proof}

\begin{theorem}[KT-THIRD-UNIT-CHARACTER-OBSTRUCTION]
Let $p>3$ be prime, let $D<0$ be a fundamental discriminant prime
to $p$, and set
\[
 L=\mathbf Q(\sqrt D,\mu_p),\qquad K=L^+,
 \qquad\theta=\chi_D\omega^{-1}.                          \tag{TU.3}
\]
Then $K/\mathbf Q$ is cyclic of degree $p-1$, and $\theta$ is a
faithful even character of its Galois group.  Moreover,
\[
 L_p(1,\theta)\equiv-B_{1,\chi_D\omega^{-2}}
       \equiv-\frac{B_{p-2,\chi_D}}{p-2}\pmod p.           \tag{TU.4}
\]
All terms are integral.  Consequently the remaining third unit of
KT-EXPLICIT-TWO-UNIT-FAMILY is equivalent to
$L_p(1,\theta)\in\mathbf Z_p^\times$.

Let $M/\mathbf Q$ be any finite multiquadratic extension and
$N=KM$.  Inflate $\theta$ to $G=\operatorname{Gal}(N/\mathbf Q)$.
With
\[
 e_\theta=\frac1{|G|}\sum_{g\in G}\theta(g)^{-1}g
          \in\mathbf Z_p[G],                              \tag{TU.5}
\]
the image of $\mathcal O_M^\times\otimes\mathbf Z_p$ in
$\mathcal O_N^\times\otimes\mathbf Z_p$ is killed by
$e_\theta$.  Thus this is also true of any collection of Pell
units from quadratic subfields, their Galois conjugates, and the
$\mathbf Z_p$-module they generate.  In particular, for $D=D_s$
in KT-EXPLICIT-TWO-UNIT-FAMILY, the explicit units used to
establish $p$-rationality of
$\mathbf Q(\sqrt{p(p+2s)})$ do not supply a nonzero vector in
the component controlling (TU.4).
\end{theorem}
\begin{proof}
The fields $\mathbf Q(\sqrt D)$ and $\mathbf Q(\mu_p)$ are
linearly disjoint: the unique quadratic subfield of the latter
is ramified at $p$, whereas the former is not.  Hence
\[
 \operatorname{Gal}(L/\mathbf Q)
       \simeq\{\pm1\}\times\mathbf F_p^\times.
\]
Under the Teichm\"uller identification, the character
$(\varepsilon,z)\mapsto\varepsilon\omega(z)^{-1}$ has kernel
$\{(1,1),(-1,-1)\}$.  The second element is complex
conjugation.  Its quotient is cyclic of order $p-1$, proving
the assertions about $K$ and $\theta$.

The character $\chi_D\omega^{-2}$ is primitive of conductor
$|D|p$, since $p>3$.  Interpolation at zero therefore gives the
exact equality
\[
                     L_p(0,\theta)=-B_{1,\chi_D\omega^{-2}}.
                                                               \tag{TU.6}
\]
There is no omitted $p$ Euler factor, since this character is
ramified at $p$.  For completeness, (TU.6) also follows directly
by substituting $s=0$ in formula~(1) on p.~2 of
Knospe--Washington.  The Bernoulli series then stops at $j=1$;
its first term is the negative defining $B_1$ sum, and its
second is one half the sum of the nontrivial character
$\chi_D\omega^{-2}$, hence zero.
Apply KT-PADIC-BRANCH-RESIDUE to $\theta$ and then the already
proved KT-FIELD-CHARACTER-REDUCTION at $n=p-2$ to obtain (TU.4).

For the projection assertion, every subextension of a
multiquadratic extension is multiquadratic.  Since $K/\mathbf Q$
is cyclic, $[K\cap M:\mathbf Q]\le2$.  Therefore
\[
 H=\operatorname{Gal}(N/M)\simeq
       \operatorname{Gal}(K/(K\cap M))
\]
has order at least $(p-1)/2>1$.  Faithfulness of $\theta$ on
$K$ implies that its restriction to $H$ is nontrivial.  Choose
$h\in H$ with $\theta(h)\ne1$.  If $u$ is the image of an
element of $\mathcal O_M^\times\otimes\mathbf Z_p$, then
$hu=u$, whereas $h(e_\theta u)=\theta(h)e_\theta u$.
The difference $\theta(h)-1$ is a $p$-adic unit, so
$e_\theta u=0$.  Also $p\nmid|G|$, so (TU.5) is integral.
The projection commutes with Galois action and all
$\mathbf Z_p$-linear operations.  This proves the claimed
scope, without an assertion about other, higher arithmetic
information that could conceivably be extracted from the same
fields.
\end{proof}

\begin{proposition}[KT-COMPOSITUM-PRATIONAL-OBSTRUCTION]
Suppose that $p>3$ is irregular, in the precise sense that for
some even $k$ with $2\le k\le p-3$ one has
$p\mid\operatorname{num}(B_k)$.  For every $D<0$ fundamental
and prime to $p$, the field
$K=(\mathbf Q(\sqrt D,\mu_p))^+$ is not $p$-rational.
Nevertheless, at $p=59$ both fields from
KT-EXPLICIT-TWO-UNIT-FAMILY satisfy the third-unit condition.
Thus full $p$-rationality of this compositum is not a necessary
condition for that third unit.
\end{proposition}
\begin{proof}
The even nontrivial character $\omega^k$ is a character of
$\mathbf Q(\mu_p)^+\subset K$.  Interpolation and (TU.1) give
\[
 L_p(1,\omega^k)\equiv L_p(1-k,\omega^k)
       =-(1-p^{k-1})\frac{B_k}{k}\equiv0\pmod p.          \tag{TU.7}
\]
Here $k$ is a unit and the ordinary Bernoulli denominator is
prime to $p$; the latter also follows from its recurrence up
to index $p-3$.  Every other nontrivial character of $K$ has
order dividing $p-1$, so its $L_p(1)$ value is integral by
KT-PADIC-BRANCH-RESIDUE.

Equation~(2.4) on pp.~4--5 of Benmerieme--Movahhedi states
that a real abelian field with character group $X$ is
$p$-rational precisely when
\[
             \prod_{1\ne\psi\in X}L_p(1,\psi)
                     \sim p^{1-m},                       \tag{TU.8}
\]
where $|\mu_{p^\infty}\cap K(\mu_p)|=p^m$ and $\sim$
means equality of valuations.  The abelian case satisfies
Leopoldt, as recalled there.  In the present case
$K(\mu_p)=L$ has degree $2(p-1)$ and contains $\mu_p$ but
not $\mu_{p^2}$, so $m=1$.  The right side of (TU.8) is a
unit, whereas the left side is integral with the nonunit
factor (TU.7).  This proves the obstruction.

For the concrete comparison, exact rational arithmetic gives
\[
 \frac{B_{44}}{44}
       =-\frac{2530297234481911294093}{2760},\qquad
       v_{59}(B_{44}/44)=1.                                \tag{TU.9}
\]
At $p=59$, the two discriminants are $D_+=-244$ and
$D_-=-228$.  The defining generalized Bernoulli sums, evaluated
modulo $59^2$ before dividing by the conductor, give
\[
 B_{1,\chi_{-244}\omega^{-2}}\equiv13,\qquad
 B_{1,\chi_{-228}\omega^{-2}}\equiv5\pmod{59}.             \tag{TU.10}
\]
The complete small certificate
\texttt{third\_unit\_59\_certificate.py} computes (TU.9)
from the Bernoulli recurrence and (TU.10) directly from the
conductor sums; it does not use the prefix formula.  The
same residues occur in the independently reviewed bounded
certificate for KT-EXPLICIT-TWO-UNIT-FINITE-CHECK.  Therefore
both third units hold, although neither corresponding $K$
is $59$-rational.
\end{proof}

\paragraph{Primary input and status.}
H.~Knospe and L.~C.~Washington, \emph{Dirichlet series
expansions of $p$-adic $L$-functions}, Abhandlungen aus dem
Mathematischen Seminar der Universit\"at Hamburg 91 (2021),
325--334, DOI \texttt{10.1007/s12188-021-00244-0}; author
version arXiv:2102.02851v3, 9 September 2021,
\url{https://arxiv.org/pdf/2102.02851}:
formula~(1) on p.~2, Proposition~2.1 and its full proof on
pp.~2--3, Theorem~2.4 and its full proof on pp.~3--5, and
Corollary~2.5 on p.~5.  These items were read in full for
the present argument.  The analytic formula imported in
that paper is Washington's Theorem~5.11; its original book
proof is not claimed freshly checked here.
For (TU.8), use the same Benmerieme--Movahhedi version and
precise pp.~4--5 locators already checked for
KT-EXPLICIT-TWO-UNIT-FAMILY.  This subsection is an
obstruction to two proposed proof mechanisms, not a new
all-prime nonvanishing or Selmer-vanishing assertion.
